\documentclass[11pt, twoside]{book}
\usepackage[letterpaper, includehead, bottom=2.5cm, left=3cm, right=3cm, top=2.5cm, headsep=1.25cm]{geometry}
\usepackage[spanish]{babel}
\usepackage{graphicx}
\graphicspath{{./images/}}
\usepackage[table]{xcolor}
\usepackage{fancyhdr}
\usepackage{lmodern}
\usepackage{enumerate, multicol}
\usepackage{enumitem}
\usepackage{tocloft}
\usepackage{amssymb}
\usepackage{float}
\usepackage{emptypage}
\usepackage[backend=biber, style=numeric]{biblatex}
    \addbibresource{bibliography.bib}

\usepackage{tikz}
\usetikzlibrary{calc, external}
%\tikzset{external/storage=./tikzfigures}
%\tikzexternalize[prefix=external-]
\usepackage{tkz-euclide}
\usetikzlibrary{patterns}

\usepackage[most]{tcolorbox}
\usepackage{imakeidx}
\usepackage{eso-pic}
\usepackage{titlesec}
\usepackage{colortbl}
\usepackage{csquotes}
\usepackage{booktabs}
\usepackage{subcaption}
\usepackage{hyperref}
%% configs and definitions

\setcounter{MaxMatrixCols}{16}


%colors
\definecolor{ustgreen}{HTML}{00685a}
\definecolor{ustmidgreen}{HTML}{009b87}
\definecolor{ustlightgreen}{HTML}{00b49d}
\definecolor{ustorange}{HTML}{ff6d3c}
\definecolor{ustdarkgreen}{HTML}{004d43}
\definecolor{graytitle}{HTML}{69827f}
\definecolor{grayheader}{HTML}{5c706d}
\definecolor{darkgraytitle}{HTML}{151a19}
\definecolor{lightbackground}{HTML}{bae7e2}
\definecolor{ustblue}{HTML}{006db6}
\definecolor{ustlightblue}{HTML}{69c0d3}
\definecolor{ustdarkblue}{HTML}{003d6b}
\definecolor{ustskyblue}{HTML}{4fa8c2}
\definecolor{usttealblue}{HTML}{006f8e}
\definecolor{ustred}{HTML}{d94c4c}
\definecolor{ustdarkred}{HTML}{9e1b32}
\definecolor{ustlightred}{HTML}{ff8074}
\definecolor{ustbrickred}{HTML}{9b3d3d}
\definecolor{ustscarlet}{HTML}{f44e42}
\definecolor{tablemid}{RGB}{232,251,247}
\definecolor{tabletop}{RGB}{192, 204, 202}

% basic definitions
\newcommand{\set}[1]{\left\lbrace #1 \right\rbrace}

%format of the title for chapters
\titleformat{\chapter}[display]
{\bfseries\large}
{\LARGE\textcolor{graytitle}{Unidad \thechapter}\\* \titlerule}
{2pt}
{\huge \centering\color{darkgraytitle}\titlerule[2pt]\vspace{0.85cm}}
[\vspace{0.85cm}{\titlerule[2pt]}]
\titlespacing*{\chapter}{0pt}{-40pt}{40pt}

% fancyhdr
\newcommand{\numberedchapter}{
    \cleardoublepage
    \fancyhead[RE]{{\bfseries \leftmark}}
    \fancyhead[LO]{{\bfseries \rightmark}}
}
\newcommand{\unnumberedchapter}[1]{
    \cleardoublepage
    \addcontentsline{t oc}{chapter}{#1}
    \fancyhead[RE]{{\bfseries #1}}
    \fancyhead[LO]{{}}
}
\pagestyle{fancy}
\fancyhf{}
\fancyhead[LE,RO]{\includegraphics[width=2.5cm]{logo_horizontal.png}}
\fancyhead[RE, LO]{\color{grayheader} Departamento de Ingeniería}
\fancyfoot[LO]{\small\color{grayheader} Dr. José Alejandro Aburto}
\fancyfoot[LE,RO]{\color{grayheader} \thepage}


% hyperref setup
\hypersetup{
    colorlinks=true,
    linkcolor=ustdarkblue,
    filecolor=black!30!magenta,
    urlcolor=black!40cyan,
    pdfauthor={Dr. José Alejandro Aburto Araneda},
}

%index
\makeindex[columns=2]

% thm & tcolorboxes numbering
\newtheorem{theorem}{Teorema}[chapter]
\newtcolorbox[auto counter, number within=chapter]{definition}[2][]{
    %float,
    colback=ustblue!2,
    colframe=ustblue,
    fonttitle=\bfseries,
    title=Definición~\thetcbcounter \hfill \textcolor{ustblue!10}{ #2 #1},
    drop fuzzy shadow=black!40!white,
}
\newtcolorbox[auto counter, number within=chapter]{example}[1][]{
    colback=ustgreen!5,
    drop fuzzy shadow=black!20!white,
    colframe=ustgreen,
    fonttitle=\bfseries,
    title=Ejemplo~\thetcbcounter,
    breakable,
    label=#1,
}
\newtcolorbox[auto counter, number within=chapter]{examples}[1][]{colback=ustgreen!5, colframe=ustgreen, fonttitle=\bfseries, #1, title=Ejemplos~\thetcbcounter, breakable}
\newtcolorbox{note}{
    enhanced,
    drop shadow=black!60!red,
    colback=ustred!7,
    frame style=
    {
        left color=ustdarkred,
        right color=ustred
    },
    fonttitle=\bfseries,
    sharp corners,
    title=¡Cuidado!,
    watermark tikz={
        \draw[line width=2mm] circle (1cm) node{\fontfamily{ptm}\fontseries{b}\fontsize{20mm}{20mm}\selectfont !};
    }
}
\newtcolorbox{simpleframe}{colframe=ustbrickred, float, enhanced}
\newtcolorbox[auto counter, number within=chapter]{thm}[2][]{
    enhanced,
    breakable,
    drop fuzzy shadow=black!50!white,
    colback=usttealblue!7,
    frame style={
        left color=usttealblue,
        right color=usttealblue!70
    },
    fonttitle=\bfseries,
    title=Teorema~\thetcbcounter \hfill \textcolor{usttealblue!13}{#1} #2
}
\newtcolorbox[auto counter, number within=chapter]{proposition}[2][]{
    enhanced,
    drop fuzzy shadow=black!40!white,
    breakable,
    colback=usttealblue!5,
    frame style={
        left color=usttealblue!80,
        right color=usttealblue!60
    },
    fonttitle=\bfseries,
    title=Proposición~\thetcbcounter \hfill \textcolor{usttealblue!13}{#1} #2
}

\newtcolorbox[auto counter, number within=chapter]{observation}[1][]{
    enhanced,
        %breakable,
        drop fuzzy shadow=black!20!white,
        colback=ustlightgreen!7,
        frame style={
            left color=ustlightgreen,
            right color=ustlightgreen!80
        },
        fonttitle=\bfseries,
        title=Observación~\thetcbcounter \hfill \textcolor{ustlightgreen!13}{#1}
    }

    \begin{document}
    \nocite{*}
    % titlepage
    \begin{titlepage}
        \newgeometry{left=7.5cm}
        \noindent
        \begin{flushright}
            \color{white}{\fontsize{32}{54} \selectfont \textsf{Matemáticas para Ciencias e Ingeniería I}}
        \end{flushright}

        \makebox[0pt][l]{\rule{1.3\textwidth}{1pt}}
        \par
        \pagecolor{ustmidgreen}

        \noindent {\Large \textcolor{ustdarkgreen}{\textsf{Universidad Santo Tomás}}\ \textcolor{ustlightgreen!10}{\textsf{Facultad de Ingeniería}}}
    \vfill
    \noindent
    \textcolor{ustlightgreen!15}
    {\large \textsf{Departamento de Ingeniería}}
    \vskip\baselineskip
    \noindent
    \textcolor{white}{\textsf{2025}}
\end{titlepage}

\restoregeometry
\nopagecolor

\pagenumbering{roman}
%\thispagestyle{empty}
\tableofcontents

\chapter*{Presentación}
\addcontentsline{toc}{chapter}{Presentación}
Este apunte busca ser una guía esencial y concisa para el curso de primer año \emph{Matemáticas para Ciencias e Ingeniería I} que comparten diversas carreras en la Universidad Santo Tomás.


Incluye el contenido en pocas palabras, incluyendo ejemplos y ejercicios propuestos. En versiones posteriores se pretende agregar:
\begin{enumerate}[label=\roman*)]
    \item Ejercicios resueltos, con mayor complejidad que los ejemplos básicos del contenido
    \item Agregar un solucionario para los ejercicios
    \item Añadir una sección o subsección en cada capitulo con ideas o problemas interesantes para profundizar
    \item Añadir una bibliografía
\end{enumerate}
\vfill
\noindent
Matemáticas para Ciencias e Ingeniería © 2025 by Dr. José Alejandro Aburto is licensed under CC BY-NC-SA 4.0. To view a copy of this license, visit https://creativecommons.org/licenses/by-nc-sa/4.0

\pagenumbering{arabic}
\chapter{Lógica y Polinomios}
\section{Lógica proposicional}
Comencemos con una definición fundamental. Queremos formalizar el concepto de \emph{proposición}, las cuales trabajaremos usualmente como incógnitas y operaremos con operaciones binarias (que requieren dos) semejantes a la suma o al producto usuales.
\begin{definition}{Proposición}
    Una proposición es una expresión que puede poseer solo uno de dos valores: \textbf{verdadero} ó \textbf{falso}.
\end{definition}
Podemos comenzar pensando en qué no es una proposición, y hay muchísimos ejemplo: todas las preguntas no son proposiciones. Primero tenemos proposiciones en el lenguaje cotidiano, y también expresiones matemáticas.
\begin{example}
    \begin{enumerate}
        \item Todos los alumnos de la UST deben sacar \(5.5\) o más para eximirse.
        \item El número \(5\) es par.
        \item La ecuación \(x^2+x+1\) no tiene solución en los números reales.
        \item Todos los gatos son grises.
    \end{enumerate}
\end{example}


Tendremos tres operaciones:
\begin{itemize}[label=\(\bullet\)]
    \item La conjunción, \enquote{y}, que denotaremos por  \(\wedge\)
    \item La disyunción, \enquote{ó}, que denotaremos por \(\vee\)
    \item La negación, \enquote{no}, que denotaremos por \(\neg\)
\end{itemize}
\begin{note}
Hay 3 notaciones para la negación, por ejemplo:
    \[\neg p=\sim p =\overline{p}\]
\end{note}
\section{Tablas de verdad}
Como las proposiciones solo pueden tener dos valores: \textbf{verdadero} ó \textbf{falso}, podemos describir los resultados de las operaciones al usarlas en una tabla. Este tipo de tabla les llamaremos \textbf{tablas de verdad}. Las siguientes tablas son por definición:
\begin{table}[h]
    \centering
    \begin{subtable}{0.3\linewidth}\centering
            \begin{tabular}{ccc}
                \toprule
                \(p\)   &   \(q\)   &   \(p\wedge q\) \\
                \midrule
                \(V\)   &   \(V\)   &   \(V\) \\
                \rowcolor{tablemid}
                \(V\)   &   \(F\)   &   \(F\) \\
                \(F\)   &   \(V\)   &   \(F\) \\
                \rowcolor{tablemid}
                \(F\)   &   \(F\)   &   \(F\) \\
            \end{tabular}
            \caption{Conjunción}
    \end{subtable}
    \hfill
         \begin{subtable}{0.3\linewidth}\centering
            \begin{tabular}{ccc}
                \toprule
                \(p\)   &   \(q\)   &   \(p\vee q\) \\
                \midrule
                \(V\)   &   \(V\)   &   \(V\) \\
                \rowcolor{tablemid}
                \(V\)   &   \(F\)   &   \(V\) \\
                \(F\)   &   \(V\)   &   \(V\) \\
                \rowcolor{tablemid}
                \(F\)   &   \(F\)   &   \(F\) \\
            \end{tabular}
            \caption{Disyunción}
    \end{subtable}
    \hfill
    \begin{subtable}{0.3\linewidth}\centering
            \begin{tabular}{cc}
                \toprule
                \(p\)   &   \(\overline{p}\) \\
                \midrule 
                \(V\)   &   \(F\)   \\
                \rowcolor{tablemid}
                \(F\)   &   \(V\)
            \end{tabular}
            \caption{Negación}
\end{subtable}
\end{table}

\begin{note}
    Las disyunción no es exclusiva, es decir, si decimos \enquote{\(p\) ó \(q\)} pueden ocurrir ambas a la vez. Hay que prestar atención a esto al principio ya que en el lenguaje cotidiano suele ser exclusivo, o sea, solemos elegir solo uno de los dos.
\end{note}


Podemos utilizar las tablas de verdad para determinar los valores posibles de una proposición compuesta, es decir, una proposición construida con estos operadores usando otras proposiciones. Por ejemplo:
\begin{example}
    Consideremos la proposición compuesta: \((p\vee q)\wedge (\overline{p}\vee\overline{q})\). Sería difícil desarrollarlo en solo un paso, por lo que agregaremos columnas a la tabla de manera que nos ayude a calcular el resultado, poniendo partes más simples que forman la proposición compuesta que nos interesa:
    \begin{center}
        \begin{tabular}{cccc}
            \toprule
            \(p\)   &   \(q\)   &    \(p\vee q\)    &  \(\overline{p}\vee\overline{q}\) \\
            \midrule
            \(V\)   &   \(V\)   &   \(V\)           & \(F\) \\
            \rowcolor{tablemid}
            \(V\)   &   \(F\)   &   \(V\)           & \(V\)\\
            \(F\)   &   \(V\)   &   \(V\)           &\(V\)\\
            \rowcolor{tablemid}
            \(F\)   &   \(F\)   &   \(F\)           & \(V\)
        \end{tabular}
    \end{center}
Ahora que ya hemos calculado estas proposiciones más simples, podemos calcular los valores posibles de la proposición compuesta inicial:
    \begin{center}
        \begin{tabular}{ccccc}
            \toprule
            \(p\)   &   \(q\)   &    \(p\vee q\)    &  \(\overline{p}\vee\overline{q}\) & \((p\vee q)\wedge (\overline{p}\vee\overline{q})\)\\
            \midrule
            \(V\)   &   \(V\)   &   \(V\)           & \(F\) & \(F\)\\
            \rowcolor{tablemid}
            \(V\)   &   \(F\)   &   \(V\)           & \(V\) &\(V\)\\
            \(F\)   &   \(V\)   &   \(V\)           &\(V\)  &\(V\)\\
            \rowcolor{tablemid}
            \(F\)   &   \(F\)   &   \(F\)           & \(V\) & \(F\)
        \end{tabular}
    \end{center}
\end{example}
Además, tenemos otro operador que aparece usualmente en matemáticas, por ejemplo en cada resultado. Solemos tener una hipótesis que es aquello que suponemos como cierto y una consecuencia de esos resultados. Podemos pensar en el teorema de Pitágoras como ejemplo, partimos de un triángulo rectángulo en el plano, y como consecuencia obtenemos una ecuación que satisfacen sus lados.
\begin{definition}{Implica}
    Definimos el conectivo (u operador)	\textbf{implica} (o \textbf{entonces}) como se indica en la siguiente tabla:
    \begin{center}
        \begin{tabular}{ccc}
            \toprule
            \(p\)   &   \(q\)   &   \(p\Rightarrow q\) \\
            \midrule
            \(V\)   &   \(V\)   &   \(V\)\\
            \rowcolor{tablemid}
            \(V\)   &   \(F\)   &   \(F\)\\
            \(F\)   &   \(V\)   &   \(V\)\\
            \rowcolor{tablemid}
            \(F\)   &   \(F\)   &   \(V\)\\
        \end{tabular}
    \end{center}
\end{definition}

Notemos que si completamos la siguiente tabla de verdad:
\begin{table}[h]
    \centering
    \begin{tabular}{cccc}
          \toprule
          \(p\)   &   \(q\)   &   \(p\Rightarrow q\)  & \(\overline{p}\vee q\)\\
            \midrule
          \(V\)   &   \(V\)   &   \(V\)               & \(V\)\\
            \rowcolor{tablemid}
          \(V\)   &   \(F\)   &   \(F\)               & \(F\)\\
          \(F\)   &   \(V\)   &   \(V\)               & \(V\)\\
            \rowcolor{tablemid}
          \(F\)   &   \(F\)   &   \(V\)               & \(V\)\\

    \end{tabular}
    \caption{Primera equivalencia lógica}
\end{table}
    Las dos últimas columnas tienen los mismos valores de verdad pero no son la misma proposición. Por ejemplo, \enquote{dos es un número par} y \enquote{5 es un número primo} tienen el mismo valor de verdad pero no son la misma proposición. Por ello, son equivalentes de cierta forma pero no iguales. Esto motiva usar un símbolo diferente:
\begin{definition}{Equivalencia lógica}
    Cuando dos proposiciones (compuestas) \(R\) y \(T\) poseen los mismos valores (en el mismo orden) en la tabla de verdad, diremos que son \textbf{equivalentes} y lo denotaremos por:
    \[R\equiv T\]
\end{definition}
Esto es una relación como equivalencia al igual que la igualdad cuando trabajamos los números reales por dar un ejemplo.

\section{Álgebra Booleana}
Ahora que tenemos un sentido de \enquote{igualdad}, quisieramos poder desarrollar expresiones como en el álgebra habitual. A esta operatoria y sus propieades les llamamos \textbf{Álgebra Booleana}.

Tenemos varias propiedades que nos serán de utilidad y enlistamos ahora:
\begin{thm}{Propiedades del álgebra booleana}
    \begin{enumerate}[label=\roman*.]
        \item Conmutatividad
            \[p\wedge q\equiv q\wedge p \qquad p\vee q \equiv q\vee q\]
        \item Identidad
            \[p\wedge V\equiv p, \quad p\wedge F\equiv F, \quad p\vee V\equiv V,\quad p\vee F\equiv p \]
        \item Idempotencia
            \[p\wedge p\equiv p, \qquad p\vee p\equiv p \]
        \item Involución
            \[\overline{\overline{p}}\equiv p\]
        \item Complemento 
            \[p\wedge \overline{p}\equiv F, \qquad p\vee \overline{p}\equiv V \]
        \item Asociatividad
            \[p\wedge (q\wedge r)\equiv (p\wedge q)\wedge r, \qquad p\vee (q\vee r)\equiv (p\vee q)\vee r \]
        \item Distributividad
            \begin{align*}
                p\wedge(q\vee r) & \equiv (p\wedge q)\vee (p\wedge r) \\
                p\vee(q\wedge r) & \equiv (p\vee q)\wedge (p\vee r)
            \end{align*}
        \item Transitividad
            \[(p\Rightarrow q)\wedge(q\Rightarrow r)\equiv p\Rightarrow r \]
        \item Leyes de De Morgan
            \[\overline{p\vee q}\equiv \overline{p}\wedge\overline{q} \qquad \overline{p\wedge q}\equiv \overline{p}\vee\overline{q}\]
        \item Absorción
            \[\left[p\wedge(p\vee q)\right]\equiv p, \qquad \left[p\vee(p\wedge q)\right]\equiv p \]
    \end{enumerate}
\end{thm}

Esto nos permite trabajar con proposiciones compuestas de una manera más convienente muchas veces.
\begin{example}
    Se tienen 3 proposiciones: \(p\), \(q\), \(r\). Sabemos que \(p\vee q \equiv r \wedge q\), queremos desarrollar y simplificar la siguiente proposición compuesta: \(\left[(p\wedge q)\vee r\right]\vee \overline{r}\). Procedemos a desarrollar utilizando las propiedades anteriores:
    \begin{align*}
        \left[(p\vee q)\wedge r\right]\vee \overline{r} & \equiv \left[(r\wedge q)\wedge r\right]\vee \overline{r} \\
                & \equiv \left[(q\wedge r)\wedge r\right]\vee \overline{r} \\
                & \equiv \left[(q\wedge (r\wedge r)\right]\vee \overline{r} \\
                & \equiv \left[q\wedge r\right]\vee\overline{r} \\
                & \equiv (q\vee\overline{r})\wedge(r\vee\overline{r}) \\
                & \equiv (q\vee\overline{r})\wedge V \\
                & \equiv q\vee\overline{r}.
    \end{align*}
\end{example}
Aunque en general podríamos obtener un resultado solo en incógnitas, a veces podremos hallar el valor.
\begin{example}
    Tenemos que
    \begin{align*}
        (p\wedge q)\vee(p\Rightarrow\overline{q})
        & \equiv (p\wedge q)\vee(\overline{p}\vee \overline{q}) \\
        & \equiv (p\wedge q)\vee \overline{\left(p\wedge q\right)} \\
        & \equiv V
    \end{align*}
\end{example}

Ya sea con álgebra booleana o con las tablas de verdad, a veces ocurrirá que sin importar los casos una proposición compuesta sea verdadera, falsa o con un valor indeterminado. Esto tiene un nombre definido como sigue:
\begin{definition}{Tautología - Contingencia - Contradicción}
    Si una proposición compuesta da como resultado
        \begin{itemize}[label=\(\bullet\)]
                \item verdadero, le llamaremos \textbf{Tautología}
                \item falso, le llamaremos \textbf{Contradicción}
                \item indeterminado, le llamaremos \textbf{Contingencia}
        \end{itemize}
\end{definition}
\begin{note}
    El que una contingencia tenga un valor indeterminado significa que depende de los valores de \(p,q,r,\ldots\).
\end{note}

\section{Funciones Proposicionales}
\begin{definition}{Función proposicional}\index{Función!proposicional}
    Una \textbf{función proposicional} es una función que toma como argumento un elemento de un universo definido y devuelve una proposición.
\end{definition}
\begin{example}
    Tenemos las siguientes funciones proposicionales:
    \begin{enumerate}[label=\roman*.]
        \item \(q(k)\equiv \text{\(k\) es par y mayor que \(5\)}\)
        \item \(F(x)\equiv x^2+x+1=0\)
        \item \(r(t)\equiv t \text{ es divisible por \(3\)}\)
    \end{enumerate}
\end{example}
En estos casos no hemos definido un universo, el cual es el dominio de la función. Sin embargo, si el contexto es claro podemos omitir esta información. 
\begin{observation}
    En general tendremos dos usos para las funciones proposicionales: Uno será con cuantificadores lógicos y el otro en la definción de conjuntos. En ambos casos estará escrito con claridad cuál es el universo en el que estaremos trabajando.
\end{observation}

\section{Cuantificadores Lógicos}
Para describir qué tanto una función proposicional nos proporciona proposiciones verdaderas, utilizaremos cuantificadores lógicos. 
\begin{definition}{Cuantificador Lógico}\index{Cuantificador lógico}
    Un \textbf{cuantificador lógico} es un operador que se aplica a las funciones proposicionales para indicar qué tantas proposiciones resultan ser verdaderas. Tenemos tres cuantificadores:
    \begin{itemize}[label=\(\bullet\)]
        \item \textbf{Para todo / Para cada}: \(\forall\)
        \item \textbf{Existe alguno / Hay algún}: \(\exists\)
        \item \textbf{Existe un único}: \(\exists!\)
    \end{itemize}
\end{definition}
\begin{example}
    \begin{enumerate}
    \item \enquote{Todos los números enteros son pares ó impares} podemos escribirlo como:
        \[\forall n\in\mathbb{Z}: n\text{ es par} \vee n\text{ es impar}\]
    \item \enquote{Existe un número par} podemos escribirlo como:
        \[\exists n\in\mathbb{Z}: n\text{ es par}\]
    \item \enquote{Para cada número natural hay un número primo mayor que él} podemos escribirlo como:
        \[\forall n\in\mathbb{N},~ \exists p\in\mathbb{N}: p\text{ es primo} ~ \wedge ~ p>n\]
    \end{enumerate}
\end{example}
Para determinar el valor de verdad de cada una de las proposiciones compuestas en las que está involucrado un cuantificador lógico, tenemos que seguir los siguientes razonamientos:
\begin{observation}
    \begin{itemize}[label=\(\bullet\)]
        \item Para que un \textbf{para todo} sea verdad, necesitamos que todos los elementos del universo nos dé verdadero al pasarlos por la función proposicional, por ello, si tan solo uno de los elementos nos da como resultado una proposición falsa, el para todo será falso.

        \item Para que un \textbf{existe} sea verdadero bastará con que uno de los elementos del universo nos de verdadero, sin importar si hay varios que nos sirvan para esto, con uno basta. Por otro lado, si todos nos dan como resultado una proposición falsa, el existe será falso.

        \item Para que un \textbf{existe un único} sea verdadero, necesitaremos que haya un elemento del universo que nos de verdadero, y además que todos los demás sean falsos. Por lo tanto, el existe un único será falso si hay dos \emph{elementos diferentes} que nos dan una proposición verdadera, ó si todos los elementos del universo nos dan falso.
    \end{itemize}
\end{observation}
\begin{example}
    \begin{enumerate}
        \item \enquote{Todos los números enteros son pares} podemos escribirlo como:
            \[\forall n\in\mathbb{Z}: n\text{ es par}\]
        En este caso, esta proposición es \textbf{falsa}. Esto, debido a que existe un\footnote{si hay más no importa, solo nos interesa que haya al menos un elemento del universo que resulte en una proposición falsa para que el \textbf{para todo} sea falso.} número entero que no es par, particularmente \(n=3\).
    \item \enquote{Existe un único número real que satisface \(x^2-1=0\)} podemos escribirlo como:
            \[\exists! x\in\mathbb{R}: x^2-1=0\]
        En este caso, esta proposición es \textbf{falsa}. Esto, debido a que existen dos números reales que satisfacen la proposición, particularmente \({x=1}\) y \({x=-1}\).
    \end{enumerate}
\end{example}
Ahora quisieramos saber como resulta ser la negación de una de estas proposiciones compuestas en las que hay cuantificadores lógicos involucrados.


Aunque ya hemos hablado acerca de cómo saber una proposición de este tipo es verdadera o falsa, no hemos visto formalmente la negación de estas proposiciones:
\begin{proposition}{Negación de Cuantificadores Lógicos}
    Tenemos las siguientes maneras de negar una proposición compuesta con cuantificadores lógicos:
    \begin{itemize}[label=\(\bullet\)]
        \item La negación del \textbf{para todo}:
            \[\neg\left(\forall u\in U: P(u)\right)\equiv \exists u\in U: \neg P(u)\]
        \item La negación del \textbf{existe}:
            \[\neg\left(\exists u\in U: P(u)\right)\equiv \forall u\in U: \neg P(u)\]
        \item La negación del \textbf{existe un único}:
            \begin{align*}
            \neg\left(\exists!\ u\in U: P(u)\right) & \equiv \left(\forall u\in U: \neg P(u)\right) \\
                                                    & \vee \exists u_1,u_2\in U: P(u_1)\wedge P(u_2)\wedge u_1\neq u_2
            \end{align*}
    \end{itemize}
\end{proposition}

\begin{example}
    Consideremos la proposición siguiente:
    \[T:\text{\enquote{El cuadrado de cada número real es mayor que 1}}\]
    La podemos escribir formamente como:
    \[T: \forall x\in\mathbb{R}: x^2>1\]
    La negación de esta proposición es:
    \[\neg T: \exists x\in\mathbb{R}: x^2\leq 1\]
    Podemos observar que \(\neg T\) es una proposición verdadera ya que por ejemplo \(x=0\) cumple que \(x^2\leq 1\). Por lo tanto, \(T\) es falsa.
\end{example}


\section{Conjuntos}
Debido a que puede resultar ser complicado definir un conjunto de una manera formal lo entenderemos como una colección de elementos. Tenemos varias maneras de representarlos:
\begin{definition}{Conjunto}
    Un \textbf{conjunto} es una colección de elementos de un universo dado. Podemos representarlo mediante:
    \begin{itemize}[label=\(\bullet\)]
        \item Definición intensiva (o por compresión): Especificamos un universo, y una propiedad (una proposición) que los elementos deben cumplir.
        \item Definición extensiva: Especificamos todos los elementos explícitamente.
        \item Definición ostensiva: Especificamos ejemplos que representan al conjunto.
    \end{itemize}
\end{definition}
\begin{note}
    Solo la definición intensiva y la extensiva son formales. Esto debido a que la definición ostensiva puede llevar a errores dependiendo de la interpretación que le demos.
\end{note}
\begin{example}[example:definiciones_conjuntos]
    \begin{enumerate}
        \item Consideremos el conjunto:
            \[\{x\in\mathbb{R}: x^2-1<0\}\]
        Este conjunto está definido de manera intensiva ó por comprensión.
        \item Consideremos el conjunto:
            \[\{2,3,5,7,9,\pi\}\]
        Este conjunto está definido de manera extensiva.
        \item Consideremos el conjunto:
            \[\{2,3,5,7,\ldots\}\]
        Este conjunto está definido de manera ostensiva.
    \end{enumerate}
\end{example}
\begin{observation}
    En el ejemplo \ref{example:definiciones_conjuntos} notemos la manera en que está escrita el conjunto
    \({x\in\mathbb{R}: x^2-1<0}\). La definición se lee como:
    \begin{center}
        \enquote{Todos los reales \(x\) tales que \(x^2-1<0\)}
    \end{center}
    Aquí, los dos puntos \enquote{\(:\)} ó bien la barra vertical \enquote{\(|\)} se lee como \enquote{tal que}.
    

    Además, el símbolo \(\in\) se lee como \enquote{pertenece a}.
\end{observation}
\begin{note}
    La manera ostensiva es solamente una representación informal. Por ejemplo el conjunto \(\{3,5,7,\ldots\}\) podemos interpretarlo de al menos dos maneras. ¿Se trata de números primos mayores que 2 o de números impares mayores que 2?
\end{note}

\section{Operaciones entre Conjuntos y Diagrama de Venn}
Tenemos tres operaciones elementales, la \textbf{unión}, la \textbf{intersección} y el \textbf{complemento}. Cada uno se obtendrá de las tres operaciones en lógica.

Denotaremos el universo por \(X\) a menos que se indique lo contrario.
\begin{definition}{Operaciones entre Conjuntos}
    \begin{itemize}[label=\(\bullet\)]
        \item \textbf{Unión}: 
            \[A\cup B=\set{x\in X: x\in A\vee x\in B}\]
        \item \textbf{Intersección}: 
            \[A\cap B=\set{x\in X: x\in A\wedge x\in B}\]

        \item \textbf{Complemento}:
            \[U\setminus A=A^c=\set{x\in X: x\notin A}\]
        \item \textbf{Diferencia}\footnote{Notemos que el símbolo que utilizamos es como un menos pero inclinado. Así lo diferenciamos del menos usual}:
            \[A\setminus B=\set{x\in X: x\in A\wedge x\notin B}\]
    \end{itemize}
\end{definition}

Para facilitar la compresión veamos unos ejemplos concretos sencillos:

\begin{example}
    Consideremos los siguientes conjuntos:
    \[A=\set{-1,1,3,6,8},\qquad B=\set{1,5,8}\]
    \[U=\set{-1,0,1,2,3,4,5,6,7,8}\]
    Y así, tenemos que:
    \begin{enumerate}[label=\alph*)]
        \item La unión de \(A\) y \(B\) es contiene los elementos de ambos conjuntos:
            \[A\cup B=\set{-1,1,3,5,6,8}\]
        \item La intersección de \(A\) y \(B\) es el conjunto de los elementos que pertenecen a ambos conjuntos a la vez:
            \[A\cap B=\set{1,8}\]
        \item El complemento de \(A\) es el conjunto de todos los elementos que no pertenecen a \(A\):
            \[U\setminus A=\set{0,2,4,5,7}\]
        \item La diferencia de \(A\) y \(B\) es el conjunto de los elementos de \(A\) que no pertenecen a \(B\):
            \[A\setminus B=\set{-1,3,6}\]
    \end{enumerate}
\end{example}
\begin{note}
    La diferencia de \(A\) y \(B\) solo quita elementos a el conjunto \(A\). Si \(B\) tiene elementos que \(A\) no posee, son totalmente ignorados. Por ejemplo, si \({A=\set{1,2}}\) y \({B=\set{2,3,4,5}}\), obtendremos que \({A\setminus B=\set{1}}\).


Notar además que \(B\setminus A=\set{3,4,5}\).
\end{note}

%% diagramas
Podemos representar los conjuntos de una manera gráfica, especialmente cuando se trata de conjuntos abstractos ó si son de finitos elementos (no son infinitos). En el caso de los conjuntos abstractos, estos diagramas nos ayudarán a comprender mejor las operaciones entre conjuntos ya que podremos visualizarlas fácilmente.

\begin{definition}{Diagrama de Venn}\index{Diagrama de Venn}
    Un \textbf{diagrama de Venn} es una representación gráfica de conjuntos que utiliza figuras geométricas cerradas para representar conjuntos. Particularmente,
         \begin{figure}[H]
                \centering
                \begin{tikzpicture}
                    \tkzDefPoints{-2.75/0/A, 2.75/0/B, 2.75/3/C, -2.75/3/D}
                    \tkzDrawPolygon(A,B,C,D)

                    \tkzDefPoint(-0.5,1.5){P1}
                    \tkzDefPoint(-0.5,0.5){P2}

                    \tkzDefPoint(0.5,1.5){Q1}
                    \tkzDefPoint(0.5,0.5){Q2}
                    % shadows
                    \tkzDrawCircle[fill=ustblue!60](P1,P2)
                    \tkzDrawCircle[fill=ustblue!60](Q1,Q2)

                    \tkzDrawCircle(P1,P2)
                    \tkzDrawCircle(Q1,Q2)


                    %labels
                \tkzDefPoint[label=right: \(B\)](1.5,1.5){L1}
                \tkzDefPoint[label=left: \(A\)](-1.5,1.5){L2}
                \tkzDefPoint[label=right: \(X\)](3,2){X}
                \end{tikzpicture}
                \caption{Unión de \(A\) con \(B\)}
        \end{figure}

         \begin{figure}[H]
                \centering
                \begin{tikzpicture}
                    \tkzDefPoints{-2.75/0/A, 2.75/0/B, 2.75/3/C, -2.75/3/D}
                    \tkzDrawPolygon(A,B,C,D)

                    \tkzDefPoint(-0.5,1.5){P1}
                    \tkzDefPoint(-0.5,0.5){P2}

                    \tkzDefPoint(0.5,1.5){Q1}
                    \tkzDefPoint(0.5,0.5){Q2}
                    % shadows
                    \begin{scope}
                        \clip (P1) circle (1);
                        \fill[ustblue!60] (Q1) circle (1);
                    \end{scope}

                    \tkzDrawCircle(P1,P2)
                    \tkzDrawCircle(Q1,Q2)

                    %labels

                    \tkzDefPoint[label=right: \(X\)](3,2){X}
                    \tkzDefPoint[label=right: \(B\)](1.5,1.5){L1}
                    \tkzDefPoint[label=left: \(A\)](-1.5,1.5){L2}
                \end{tikzpicture}
                \caption{Intersección de \(A\) con \(B\)}
        \end{figure}
         \begin{figure}[H]
                \centering
                \begin{tikzpicture}
                    \tkzDefPoints{-2.75/0/A, 2.75/0/B, 2.75/3/C, -2.75/3/D}
                    \tkzDrawPolygon(A,B,C,D)

                    \tkzDefPoint(-0.5,1.5){P1}
                    \tkzDefPoint(-0.5,0.5){P2}

                    \tkzDefPoint(0.5,1.5){Q1}
                    \tkzDefPoint(0.5,0.5){Q2}

                    \tkzDrawCircle[fill=ustblue!60](P1,P2)
                    \tkzDrawCircle[fill=white](Q1,Q2)
                    \tkzDrawCircle(P1,P2)

                    %labels

                    \tkzDefPoint[label=right: \(X\)](3,2){X}
                    \tkzDefPoint[label=right: \(B\)](1.5,1.5){L1}
                    \tkzDefPoint[label=left: \(A\)](-1.5,1.5){L2}
                \end{tikzpicture}
                \caption{Diferencia de \(A\) con \(B\)}
        \end{figure}
         \begin{figure}[H]
                \centering
                \begin{tikzpicture}
                    \tkzDefPoints{-2.75/0/A, 2.75/0/B, 2.75/3/C, -2.75/3/D}
                    \tkzDrawPolygon[fill=ustblue!60](A,B,C,D)
                    \tkzDrawPolygon(A,B,C,D)

                    \tkzDefPoint(0,1.5){P1}
                    \tkzDefPoint(0,0.5){P2}


                    \tkzDrawCircle[fill=white](P1,P2)
                    \tkzDrawCircle(P1,P2)

                    \tkzDefPoint[label=right: \(X\)](3,2){X}
                    \tkzDefPoint[label=right: \(A\)](1.5,1.5){L1}
                \end{tikzpicture}
                \caption{Complemento de \(A\)} 
        \end{figure}
\end{definition}
Podemos representar diferentes operaciones entre conjuntos usando el diagrama de Venn.
\begin{example}
        Consideremos el siguiente diagrama de Venn que representa una operación entre los conjuntos \(A\) y \(B\):
             \begin{figure}[H]
                \centering
                \begin{tikzpicture}
                    \tkzDefPoints{-2.75/0/A, 2.75/0/B, 2.75/3/C, -2.75/3/D}
                    \tkzDrawPolygon(A,B,C,D)

                    \tkzDefPoint(-0.5,1.5){P1}
                    \tkzDefPoint(-0.5,0.5){P2}

                    \tkzDefPoint(0.5,1.5){Q1}
                    \tkzDefPoint(0.5,0.5){Q2}
                    % shadows
                    \tkzDrawCircle[fill=ustblue!60](P1,P2)
                    \tkzDrawCircle[fill=ustblue!60](Q1,Q2)
                    \begin{scope}
                        \clip (P1) circle (1);
                        \fill[white] (Q1) circle (1);
                    \end{scope}

                    \tkzDrawCircle(P1,P2)
                    \tkzDrawCircle(Q1,Q2)

                    %labels

                    \tkzDefPoint[label=right: \(X\)](3,2){X}
                    \tkzDefPoint[label=right: \(B\)](1.5,1.5){L1}
                    \tkzDefPoint[label=left: \(A\)](-1.5,1.5){L2}
                \end{tikzpicture}
                \caption{Diferencia Simétrica de \(A\) con \(B\)}
        \end{figure}
        En esta figura tenemos representado lo que es conocido como \textbf{diferencia simétrica}. Gracias a la figura, podemos deducir que:
        \[A\setminus B \cup B\setminus A = \left(A\cup B\right)\setminus\left(A\cap B\right)\]
\end{example}
\section{Cardinalidad}
    Cuando tenemos un conjunto, sabemos si un elemento pertenece al conjunto o no.
    Considerando aquellos que sí pertenecen podemos asignarle un valor, y este valor será llamado cardinalidad.
\begin{definition}{Cardinalidad}\index{Cardinalidad}
    La cardinalidad de un conjunto es el número de elementos que contiene. Sea el conjunto \(A\), denotaremos su cardinalidad
por: \(|A|\).
\end{definition}
\begin{note}
    Es importante que solo contemos una vez cada elemento del conjunto, ya que por diversas circunstancias podría aparecer varias veces.
\end{note}

\begin{example}
    Calculemos la cardinalidad de los siguientes conjuntos:
        \begin{enumerate}[label=\alph*)]
            \item \(\left|\set{-3,4,5,6,7}\right|=5\)
            \item \(\left|\set{\frac{1}{2},\frac{2}{4},\frac{6}{12},-\frac{1}{2}}\right|=2\)
            \item \(\left|\set{2,3,5,7,11,13}\right|=6\)
        \end{enumerate}
\end{example}

% Esto hay que explicarlo mejor y añadir diagramas de Venn para mejorar la compresión
Cuando se trata conjuntos con cardinalidad finita, podemos obtener la cardinalidad de un conjunto formado mediante operaciones de varios conjuntos vía dos propiedades principalmente:
\begin{proposition}{Propiedades de la Cardinalidad}
    Sean \(A\) y \(B\) dos conjuntos. Entonces:
        \[|A\cup B|=|A|+|B|-|A\cap B\]
        \[|A\setminus B|=|A|-|A\cap B|\]
\end{proposition}
\begin{observation}
    Vale la pena notar que si dos conjuntos son disjuntos, es decir, que su interesección es vacía. Entonces, su cardinalidad es la suma de las cardinalidades de ambos conjuntos.
\end{observation}

\begin{example}
    Si \(|A|=3\) y \(|B|=4\) y además \(|A\cup B|=5\), y queremos determinar \(|A\cap B|\), entonces tenemos que
    \begin{align*}
        |A\cup B| & = |A|+|B|-|A\cap B| \\
        5 & = 3+4-|A\cap B| \\
        -|A\cap B|      & = 5-3-4 \\
         |A\cap B|      & = 2
    \end{align*}
\end{example}

\begin{example}
    En una clase, cada estudiante debe elegir algún taller para realizar. De los 40 estudiantes,
    15 de ellos eligieron el taller de arduino, 20 de ellos el taller de ajedrez y 15 de ellos el taller de juego de mesa.
    Se sabe que:
    \begin{itemize}[label=\(\bullet\)]
        \item No hay estudiantes que tengan el taller de ajedrez y el de juegos de mesa.
        \item Hay 5 estudiantes que solo tienen el taller de arduino.
        \item Hay 7 estudiantes que tienen el taller de arduino y el de ajedrez.
    \end{itemize}
    ¿Cuántos estudiantes tienen el taller de juegos de mesa y Arduino juntos? ¿Cuántos solo juegos de mesa?
\end{example}
\section{Conjuntos numéricos}
??

\section{Polinomios}
\begin{definition}{Polinomio}\index{Polinomio}
    Un polinomio es una expresión algebraica de la forma siguiente:
        \[a_nx^n+a_{n-1}x^{n-1}+\cdots+a_1x+a_0\]
    en donde \(a_j\) son coeficientes de algún conjunto en particular, y \(x\) es una variable.

    Denotaremos al conjunto de los polinomios de variable \(x\) y coeficientes en un conjunto \(A\) por \(A[x]\).
\end{definition}
En este caso los coeficientes pertenecerán a un conjunto numérico, por ejemplo \(\mathbb{Z}\).

También tenemos que conocer los nombres usuales que utilizaremos para referirnos a ciertas características o partes de un polinomio.

%valdra la pena hablar de termino?? casi nunca se usa eso
\begin{definition}{Grado, coeficiente libre y lider}
    Sea \(p(x)\in A[x]\), digamos \(p(x)=a_nx^n+a_{n-1}x^{n-1}+\cdots+a_1x+a_0\). Entonces:
        \begin{itemize}[label=\(\bullet\)]
            \item El grado de \(p(x)\) es el mayor exponente de \(x\) que aparece en \(p(x)\).
            \item El coeficiente libre de \(p(x)\) es el coeficiente libre de la variable \(x\) en \(p(x)\).
            \item El lider de \(p(x)\) es el coeficiente del t\'ermino de mayor grado en \(p(x)\).
        \end{itemize}
\end{definition}
Veamos algunos ejemplos para clarificar esta definici\'on:
\begin{example}
    \begin{enumerate}
        \item Consideremos el polinomio: \(p(x)=x^3-2x^2+3x-1\). Entonces,
            \begin{itemize}[label=\(\bullet\)]
                \item El grado de \(p(x)\) es \(3\).
                \item El coeficiente libre de \(p(x)\) es \(-1\).
                \item El lider de \(p(x)\) es \(1\).
            \end{itemize}
        \item Consideremos el polinomio: \(q(x)=4x^3+2x\). Entonces,
            \begin{itemize}[label=\(\bullet\)]
                \item El grado de \(q(x)\) es \(3\).
                \item El coeficiente libre de \(q(x)\) es \(2\).
                \item El lider de \(q(x)\) es \(4\).
            \end{itemize}
        \item Consideremos el polinomio: \(r(x)=-x^2-x+6\). Entonces,
            \begin{itemize}[label=\(\bullet\)]
                \item El grado de \(r(x)\) es \(2\).
                \item El coeficiente libre de \(r(x)\) es \(6\).
                \item El lider de \(r(x)\) es \(-1\).
            \end{itemize}
    \end{enumerate}
\end{example}
\begin{observation}
    El grado de un polinomio \(p(x)\) lo denotaremos por \(\deg(p(x))\). Además, notemos que el grado del polinomio nulo no quedaría bien definido si le asignamos un número entero. 

    Convendremos que \(\deg(0)=-\infty\).
\end{observation}

%% mucho texto
En los polinomios tenemos dos operaciones, que son la suma y el producto. Los cuales hereder\'an sus propiedades de la suma y producto que ya tenemos en el conjunto \(A\). 
% valdra la pena profundizar en las operaciones de los polinomios?? igual hay gente que podria no saber distribuir o no trabajar apropiadamente con el algebra.

Cuando se trata realizar operaciones entre polinomios, hay algunas expresiones que aparecen regularmente y por ello vale la pena conocer, e incluso memorizar ciertas reglas para poder trabajar de una manera más c\'omoda. Estas expresiones recurrentes son conocidas como \textbf{productos notables}, debido a que justamente se trata de ciertos productos entre polinomios.


Tendremos una herramienta simple para utilizaremos en algunas potencias de binomios (polinomios de dos t\'erminos).
\begin{definition}{Triángulo de Pascal}\index{Triángulo!de Pascal}
    Sea \(n\in\mathbb{N}\). El tri\'angulo de Pascal de \(n\) filas es el siguiente arreglo de n\'umeros:
        \[\begin{matrix}
            & & & & & 1 & & & & & \\
            & & & & 1 & & 1 & & & &\\
            & & & 1 & & 2 & & 1 & & & \\
            & & 1 & & 3 & & 3 & & 1 & &\\
            & 1 & & 4 & & 6 & & 4 & & 1 & \\
            \reflectbox{\(\ddots\)} &&&&&\vdots &&&&& \ddots
            \end{matrix}\]
        En los extremos ponemos 1, y rellenamos con la suma de los números de la fila anterior.
\end{definition}

Utilizaremos el triángulo de Pascal para poder determinar los coeficientes de las potencias de un binomio, es decir, la potencia de una suma de dos términos.

\begin{example}
    Escribimos los desarrollos de las siguientes poencia de binomio utilizando los coeficientes del triángulo de Pascal:
    \begin{itemize}[label=\(\bullet\)]
        \item \((x+y)^0=1\)
        \item \((x+y)^1=x+y\)
        \item \((x+y)^2=x^2+2xy+y^2\)
        \item \((x+y)^3=x^3+3xy^2+3xy+y^3\)
        \item \((x+y)^4=x^4+4xy^3+6xy^2+y^4\)
    \end{itemize}
\end{example}
Podemos seguir escribiendolos ampliando el triángulo de Pascal.


Vale la pena igualmente conocer otros productos notables que no aparecen naturalmente de una potencia como las anteriores:
%% sumas por su diferencia y esa clase de cosas 


\section{División de polinomios}
Al igual que en la división de números enteros, podemos realizar una división de polinomios.


Primero que todo, esta división se corresponde con el siguiente resultado:

\begin{thm}{Algoritmo de Euclides}
    Sean \(f(x)\in A[x]\) y \(g(x)\in A[x]\). Entonces, existen dos polinomios \(h(x), r(x)\in A[x]\) tales que:
        \[f(x)=g(x)h(x)+r(x),\]
    en donde \(r(x)=0\) ó \(\deg(r(x))<\deg(q(x))\).

    A \(r(x)\) se le llama el \textbf{resto} de la división de \(f(x)\) por \(g(x)\). Además, diremos que \(f(x)\) es el dividendo, \(g(x)\) es el divisor y \(h(x)\) es el cuociente.
\end{thm}

Este resultado puede resultar se puramente abstracto, y aún nos queda ver cómo se realiza la división de polinomios.
\begin{observation}
    Utilizaremos los siguientes polinomios para ejemplificar el algoritmo de Euclides: \(f(x)=4x+x^2+1\) y \(g(x)=x+1\).
    Para dividir \(f(x)\) por \(g(x)\) aplicaremos los siguientes pasos:
        \begin{itemize}[label=\(\bullet\)]
            \item Paso 1: Primero ponemos los polinomios ordenando sus términos de mayor a menor grado, de izquierda a derecha:
                \[x^2+4x+1:x+1\]
            \item Paso 2: Buscamos un monomio que al multiplicarlo por \(x+1\) resulte en un polinomio cuyo término lider coincida con el de \(x^2+4x+1\).
Ese monomio siempre es único. Este este caso, es \(x\), ya que \(x(x+1)=x^2+x\), que tiene el mismo término lider que \(x^2+4x+1\). Lo anotamos de la manera siguiente:
\[x^2+4x+1:x+1={\color{ustdarkblue}x}\]
            \item Paso 3: Anotamos el producto de \(x\) con \(x+1\) abajo de la división:
                \begin{align*}
                    & x^2+4x+1: x+1 = x\\
                    - & {\color{ustdarkblue}(x^2+x)}
                \end{align*}
            \item Paso 4: Realizamos la resta del dividendo menos el producto recién mencionado:
                \begin{align*}
                    & {\color{ustdarkblue}x^2+4x+1}: x+1 = x\\
                    - & {\color{ustdarkblue}(x^2+x)}\\
                    = & ~{\color{ustdarkred}3x+1}
                \end{align*}
            \item Paso 5: Repetimos el proceso con el nuevo dividendo \(3x+1\), y repetimos hasta que el resultado de la resta tenga grado menor que el divisor, o que sea cero. En este caso:
                 \begin{align*}
                    & x^2+4x+1: x+1 = x+{\color{ustdarkblue}3}\\
                    - & (x^2+x)\\
                    = & ~3x+1 \\
                    \color{ustdarkblue}- & {\color{ustdarkblue}(3x+3)}\\
                = & ~{\color{ustdarkred}-2}
                \end{align*}
        \end{itemize}
    La conclusión de este algoritmo es que:
        \[x^2+4x+1=(x+1)(x+3)-2.\]
\end{observation} 

\begin{example}
    Dividimos \(f(x)=x^3+2x-3\) por \(g(x)=x-1\), y resulta: 
        \[x^3+2x-3=(x-1)(x^2+x+3).\]
    En este caso el resto es cero.
\end{example}
\begin{example}
 Dividimos \(f(x)=x^3+1\) por \(g(x)=x^2\),  y resulta:
        \[x^3+1=x\cdot x^2 +1.\]
        Aquí el resto es \(1\).
\end{example}

\section{Teorema del Resto}
\begin{definition}{Ra\'iz de un polinomio}\index{Ra\'iz!de un polinomio}
    Sea \(p(x)\in A[x]\). Una ra\'iz de \(p(x)\) es el elemento \(a\in A\) tal que \(p(a)=0\).
\end{definition}
Ocurrirá que cada vez que un polinomio tenga una raíz, se podrá factorizar en un producto de polinomios en donde estará presente la raíz.
%% valdra la pena poner un ejemplo motivador antes??
\begin{thm}{Teorema del Resto}
    Sea \(p(x)\in A[x]\) y \(a\in A\). Entonces:
        \[p(x)=(x-a)q(x)+p(a).\]
En otras palabras, si dividimos \(p(x)\) por \((x-a)\) el resto resultará ser \(p(a)\).
\end{thm}

Notemos que en particular, si \(a\) es una raíz de \(p(x)\), entonces \(p(a)=0\).
\begin{thm}{Teorema del factor}
    Sea \(p(x)\in A[x]\) y \(a\in A\) es una raíz de \(p(x)\). Entonces existe un polinomio \(q(x)\in A[x]\) tal que
        \[p(x)=(x-a)q(x).\]
\end{thm}

\section{Raíces Reales y Racionales}
Cuando queremos hallar las raíces de un polinomio de grado 2, podremos encontrarlas utilizando la fórmula de Bhaskara. Sin embargo, cuando se trata de otros polinomios de mayor grado, la dificultad crece rápidamente. Hace siglos, la primera esperanza era hallar una solución por radicales, como en la fórmula de Bhaskara o la de Cardano, pero tomo varios siglos en descubrir que es imposible en general para polinomios de grado 5 o más. 

Hay dos maneras de enfrentar este problema: Hallar las soluciones más simples de encontrar en polinomios amigables, como los que tienen coeficientes enteros. Y la otra alternativa es aproximar esas soluciones. Nosotros aprenderemos la primera de estas maneras en esta sección.
\begin{thm}{Teorema de Gauss}
    Sea \(p(x)\in\mathbb{Z}\), digamos \(p(x)=a_nx^n+a_{n-1}x^{n-1}+\cdots+a_1x+a_0\), con \(a_n\neq 0\). Entonces, las raíces racionales de \(p(x)\) pertenecen al conjunto:
    \[\set{\pm\frac{c}{d}\;:\; c \text{ es un divisor de } a_0 \wedge d \text{ es un divisor de } a_n}\]
\end{thm}
Primero tenemos que entender como se contruye este conjunto, a pesar de que en su definción se describe completamente, es buena idea aclarar los conceptos ante ideas nuevas.
\begin{example}
    Consideremos el polinomio \(p(x)=x^3-3x^2-8x+4\). Entonces, el conjunto de las posibles raíces racionales de \(p(x)\) es:
        \[\set{\pm\frac11,\pm\frac21,\pm\frac41}\]
    Y probamos estos valores evaluando en \(p(x)\) uno a uno:
        \begin{align*}
            p\left(1\right)&=1^3-3\cdot1^2-8\cdot1+4 = -6\\
            p\left(2\right)&=2^3-3\cdot2^2-8\cdot2+4 = -16\\
            p\left(4\right)&=4^3-3\cdot4^2-8\cdot4+4 = -12\\
            p\left(-1\right)&=(-1)^3-3\cdot(-1)^2-8\cdot(-1)+4 =8 \\
            p\left(-2\right)&=(-2)^3-3\cdot(-2)^2-8\cdot(-2)+4 =0 \\
            p\left(-4\right)&=(-4)^3-3\cdot(-4)^2-8\cdot(-4)+4 = -76
        \end{align*}
        Y hemos encontrado una raíz racional de \(p(x)\).
\end{example}

\section{Números complejos}
Los números complejos se documentaron por primera vez en el siglo XVI cuando Cardano encontró una solución en radicales para la cúbica. Desde entonces ha sido objeto de estudio debido a su utilización en Electromagnetismo, Mecánica cuántica, acústica, electroencefalografía, resonancias magnéticas, etc.

Ahora nos corresponde estudiar las bases de este interesante objeto matemático. 


\begin{definition}{Conjunto de números complejos}\index{Números complejos}
El conjunto de los números complejos se define por:
    \[\mathbb{C}=\set{a+bi: a,b\in\mathbb{R}}\]
    en donde \(i^2=-1\).
\end{definition}

Las operaciones entre números complejos estás definidas polinomialmente. Aunque esto ya lo hemos visto, es importante nota que la adición de la unidad imaginaria \(i\) propocionará diferencias notables.
\begin{proposition}{Operaciones en \(\mathbb{C}\)}
    Las operaciones algebraicas en \(\mathbb{C}\) son la suma, el producto, la conjugación y la inversión. Sean \(a,b,c,d\in\mathbb{R}\), luego se tiene que:
    \begin{enumerate}[label=\roman*)]
        \item La suma:
            \[\left(a+bi\right)+\left(c+di\right)=\left(a+c\right)+\left(b+d\right)i\]
        \item El producto:
            \[\left(a+bi\right)\left(c+di\right)=\left(ac-bd\right)+\left(ad+bc\right)i\]
        \item La conjugación:
        \[\overline{\left(a+bi\right)}=\left(a-bi\right)\]
        \item La inversión:
            \[\frac{1}{a+bi}=\]
    \end{enumerate}
\end{proposition}

\newpage
\section{Ejercicios}
\subsection*{Lógica}
\begin{enumerate}
    \item Calcule la tabla de verdad de la siguiente proposición:
            \[\left[p\vee(\overline{p}\wedge q)\right]\Rightarrow q\]
    \item Usando tablas de verdad, pruebe que la siguiente proposición es una tautología:
        \[\left(p\Rightarrow q\right) \iff \left(\overline{q}\Rightarrow \overline{p}\right)\]
    \item Usando tablas de verdad, pruebe que la siguiente proposición es una contingencia:
        \[(p\vee q) \Rightarrow (\overline{p}\wedge \overline{q})\] 
\end{enumerate}
\subsection*{Conjuntos}
\begin{enumerate}
    \item Sean los conjuntos \(A=\set{1,3,5,7}\), \(B=\set{1,3,5,9}\), \(C=\set{9}\).
        Determine si son verdaderas o falsas las siguiente contenciones:
        \begin{multicols}{2}
        \begin{enumerate}[label=\alph*)]
            \item \(A\subset B\)
            \item \(B\subset A\)
            \item \(A\subset C\)
            \item \(B\subset C\)
            \item \(C\subset A\)
            \item \(C\subset B\)
        \end{enumerate}
        \end{multicols}
    \item Determine la cardinalidad de los siguientes conjuntos:
        \begin{enumerate}[label=\alph*)]
            \item \(\set{-1, 1,2}\)
            \item \(\set{\frac39,\frac13,\frac19}\)
            \item \(\set{2,3,5,7,-\frac42,\frac93}\)
        \end{enumerate}
\end{enumerate}
\subsection*{Polinomios}
\begin{enumerate}
    \item 
    \item
    \item Factorice y simplifique la siguiente expresión:
    \item 
\end{enumerate}
\subsection*{Números complejos}


\chapter{Geometría Analítica}
\section{El plano Euclideano}


\chapter{Trigonometría y Relación de Orden}
\section{Trigonometría en el Triángulo}
\section{Trigonometría en el plano}
\section{Identidades Trigonometricas}

\section{Forma polar de un número complejo}



\chapter{Funciones y Límites}
\section{Funciones reales}
\section{Funciones polinomiales, racionales y raíces}
\section{Funciones Trigonométricas y sus inversas}
\section{Función exponencial y logaritmos}
\section{Límites y Continuidad}
\section{Álgebra de Límites}
\section{Cambios de variable}
\section{Ejercicios}




\printbibliography%
\printindex

\begin{titlepage}
    \newgeometry{left=7.5cm}

    \makebox[0pt][l]{\rule{1.3\textwidth}{0pt}}
    \par
    \pagecolor{ustmidgreen}

    \noindent
\end{titlepage}

\nopagecolor
\end{document}
